\chapter{部分積分}\label{章:部分積分}

前章では，合成関数の微分を逆向きに読んで，複雑なまとまりを一つの変数に換えた。たとえば $x(x^2+1)^{20}$ には，$x^2+1$ の微分に当たる $2x$ が，定数倍を除いて含まれていた。ところが $x\e^x$ では，$x$ と $\e^x$ のどちらかを新しい変数にしても，もう片方と $\dd x$ が一緒に簡単な形へ変わるわけではない。

ここで，積の微分を思い出す。$x\e^x$ を微分すると，$x\e^x$ のほかに $\e^x$ が現れる。したがって，$x\e^x$ 自体は求める原始関数ではないが，余分に現れる $\e^x$ の原始関数を引けばよさそうである。このように，積を微分したときの二つの項を使って，ある積分を別の積分へ換えるのが部分積分である。

\section{積の微分から公式を作る}\label{節:部分積分の導出}

二つの関数 $u(x),v(x)$ が微分できるとき，
\[
 \{u(x)v(x)\}'=u'(x)v(x)+u(x)v'(x)
\]
である。右辺の第二項を積分したいなら，この等式を
\[
 u(x)v'(x)=\{u(x)v(x)\}'-u'(x)v(x)
\]
と書き直す。両辺の原始関数を考えると，積 $uv$ と，もう一つの積分が現れる。

\begin{theorem}{部分積分の公式}
$u,v$ がある区間で連続微分可能ならば，その区間で
\[
 \int u(x)v'(x)\dd x
 =u(x)v(x)-\int u'(x)v(x)\dd x
\]
が成り立つ。不定積分は任意定数を含めた原始関数の族として読む。特に，同じ区間の $a,b$ に対して
\[
 \int_a^b u(x)v'(x)\dd x
 =[u(x)v(x)]_a^b-\int_a^b u'(x)v(x)\dd x
\]
が成り立つ。
\end{theorem}

定積分の公式では，積の微分を $a$ から $b$ まで積分し，基本定理を使った。$[uv]_a^b$ は $u(b)v(b)-u(a)v(a)$ である。この端点の項を省くことはできない。積の微分のうち，一方の項だけを移項したことに対応する量である。

公式の左辺では $u$ をそのまま，$v'$ を微分された形で書いている。右辺へ移ると，$u$ は $u'$ へ微分され，$v'$ は原始関数 $v$ に変わる。したがって実際の選択は，「どの因子を微分し，どの因子を積分すると，残る積 $u'v$ が扱いやすくなるか」という問題になる。

ときには $\dd v=v'(x)\dd x$，$\dd u=u'(x)\dd x$ とまとめて，
\[
 \int u\dd v=uv-\int v\dd u
\]
とも書く。これは同じ公式の短い表記であり，$u$ と $v$ を関係のない変数として自由に動かしているのではない。本章では，微分する因子と積分する因子を確かめやすいように，主に $u,v'$ を書いて計算する。

\begin{example}{一次式を微分して定数にする}
$\int x\e^x\dd x$ を求めよ。
\end{example}

$x$ は微分すると $1$ になる。一方，$\e^x$ は積分しても $\e^x$ なので，この二つを組み合わせれば残る積分が簡単になる。そこで
\[
 u=x,\qquad v'=\e^x,
 \qquad u'=1,\qquad v=\e^x
\]
と選ぶ。部分積分から
\[
\begin{aligned}
 \int x\e^x\dd x
 &=x\e^x-\int 1\cdot\e^x\dd x\\
 &=x\e^x-\e^x+C
 =(x-1)\e^x+C.
\end{aligned}
\]
積 $x\e^x$ を作って終わりにすると，その微分は $x\e^x+\e^x$ となる。余分な $\e^x$ を取り除くために $\e^x$ を引いている。答えの微分も
\[
 \frac{\dd}{\dd x}\{(x-1)\e^x\}
 =\e^x+(x-1)\e^x=x\e^x
\]
となり，この考え方をそのまま確認している。

逆に $u=\e^x$，$v'=x$ と選ぶことも，公式上は正しい。その場合には
\[
 \int x\e^x\dd x
 =\frac{x^2}{2}\e^x-\frac12\int x^2\e^x\dd x
\]
となる。ところが，もとの一次式が二次式になり，残った積分はむしろ手間が増えている。部分積分の公式を使えることと，その選択で計算が進むことは別である。右辺に何が残るかを一度書いてから，因子の選び方を判断すればよい。

\section{微分すると簡単になる因子を選ぶ}\label{節:部分積分の因子選択}

整式は微分すると次数が下がるため，微分する側の候補になる。指数関数や三角関数は，積分しても同じ種類の関数のままなので，積分する側に選びやすい。ただし，この順序は絶対の規則ではない。対数を含む積では，対数を微分すると分数になり，別の因子と約分できることがある。

また，見た目が積でなくても，$1$ を掛けた積と考えられる。$\log x$ の原始関数を直接思い出せない場合でも，$\log x\cdot1$ と見れば，積分する側に既知の関数 $1$ を置ける。

\begin{example}{一つの関数を積と見る}
$x>0$ において，$\int\log x\dd x$ を求めよ。
\end{example}

$\log x$ の微分は $\frac1x$ である。これに $1$ の原始関数 $x$ が掛かれば，積は $1$ になる。そこで
\[
 u=\log x,\qquad v'=1,
 \qquad u'=\frac1x,\qquad v=x
\]
と選ぶと，
\[
\begin{aligned}
 \int\log x\dd x
 &=x\log x-\int x\cdot\frac1x\dd x\\
 &=x\log x-x+C.
\end{aligned}
\]
もとの対数を，微分して得られる分数と，積分して得られる $x$ の組合せで消したのである。微分による確認は
\[
 \frac{\dd}{\dd x}(x\log x-x)
 =\log x+x\cdot\frac1x-1=\log x
\]
となる。$\log x$ は $x>0$ で考えているので，この計算にも同じ条件が付く。

この選び方は $x\log x$ にも使える。$\log x$ を微分し，$x$ を積分すると，
\[
\begin{aligned}
 \int x\log x\dd x
 &=\frac{x^2}{2}\log x
   -\int\frac{x^2}{2}\cdot\frac1x\dd x\\
 &=\frac{x^2}{2}\log x-\frac{x^2}{4}+C.
\end{aligned}
\]
ここでは整式 $x$ を積分する側に選んだ。それでも，対数を微分して得た $\frac1x$ と掛けることで，残る積分が単なる一次式になる。関数の名前だけで順序を決めず，$u'v$ を具体的に見ることが大切である。

一方，$\int\frac{\log x}{x}\dd x$ ならば，前章で見たとおり $t=\log x$ と置くほうが短い。$\frac1x\dd x=\dd t$ となり，ただちに $\int t\dd t$ へ変わるからである。「対数があるから部分積分」と決めるのではなく，対数の微分がすでに残りの部分にあるかを先に調べよう。

因子の選択には，もう一つ条件がある。積分する側の原始関数 $v$ を，その時点で求められなければならない。たとえば $\int x\e^{x^2}\dd x$ で，$x$ を微分しようとして $v'=\e^{x^2}$ と選ぶと，そこで原始関数が分からなくなる。しかし $t=x^2$ と置けば，$x\dd x=\frac12\dd t$ なので，答えは $\frac12\e^{x^2}+C$ である。部分積分を試す前に，置換に必要な微分の因子がそろっていないかを見ると，遠回りを避けられる。

\section{同じ選択を繰り返す}\label{節:部分積分の反復}

整式と指数関数，あるいは整式と三角関数の積では，一度の部分積分で整式の次数が一つ下がる。残る積分に同じ方針を適用すれば，やがて整式が定数になる。最初から一度で計算を終わらせようとする必要はない。何回か繰り返せば，確実に既知の積分へ到達することが見えている。

\begin{example}{二次式の次数を一つずつ下げる}
$\int x^2\cos x\dd x$ を求めよ。
\end{example}

まず $u=x^2$，$v'=\cos x$ と選ぶ。$u'=2x$，$v=\sin x$ なので，
\[
 \int x^2\cos x\dd x
 =x^2\sin x-2\int x\sin x\dd x.
\]
残った積分では，$x$ を微分し，$\sin x$ を積分する。$\sin x$ の原始関数は $-\cos x$ だから，
\[
\begin{aligned}
 \int x\sin x\dd x
 &=x(-\cos x)-\int1\cdot(-\cos x)\dd x\\
 &=-x\cos x+\sin x+C.
\end{aligned}
\]
これをもとの式へ戻して，任意定数をまとめると
\[
\begin{aligned}
 \int x^2\cos x\dd x
 &=x^2\sin x-2(-x\cos x+\sin x)+C\\
 &=x^2\sin x+2x\cos x-2\sin x+C.
\end{aligned}
\]
二回目の積分で現れた負号だけでなく，その式全体に最初の $-2$ が掛かることも確かめる。途中の積分を独立に計算したときは，戻す先の係数を残しておくと，符号の取り違えを防げる。

答えを微分すると，途中で現れた項が対になって消える。
\[
\begin{aligned}
 &\frac{\dd}{\dd x}
  \{x^2\sin x+2x\cos x-2\sin x\}\\
 &\qquad=2x\sin x+x^2\cos x
       +2\cos x-2x\sin x-2\cos x\\
 &\qquad=x^2\cos x.
\end{aligned}
\]
部分積分で補った項は，このように一つ前の微分で生まれた余分な項を打ち消している。繰り返す計算は，思い付きで項を増やす操作ではなく，積の微分による余分な項を順に取り除く操作である。

一般に，$P(x)$ が $n$ 次の整式なら，$P,P',P'',\ldots$ と微分を繰り返すと，$n$ 回目に定数になり，その次には $0$ になる。したがって $P(x)\e^{ax}$ や $P(x)\sin(ax)$，$P(x)\cos(ax)$ は，$a\ne0$ のとき，整式を微分する部分積分を繰り返して計算できる。$a$ が付く場合は，指数関数や三角関数を積分するたびに $a$ で割る。次数が下がることと，係数が変わることを別々に確認すればよい。

\section{もとの積分が戻ったら等式を解く}\label{節:部分積分の循環}

指数関数と三角関数の積では，整式のように微分を繰り返して $0$ になる因子がない。その代わり，正弦と余弦は微分によって互いに移り変わる。二回の部分積分で，もとの積分が再び現れることがある。そこで計算を失敗と決めず，もとの積分を未知数として含む等式が得られていないかを調べる。

\begin{example}{指数関数と正弦の積}
$\int\e^x\sin x\dd x$ を求めよ。
\end{example}

積分定数に惑わされずに等式を整理するため，まず
\[
 I(x)=\int_0^x\e^t\sin t\dd t,\qquad
 J(x)=\int_0^x\e^t\cos t\dd t
\]
と定める。$I,J$ は $x$ の関数であり，基本定理から $I'(x)=\e^x\sin x$ である。したがって $I(x)$ が求まれば，それに任意定数を加えて不定積分の答えにできる。

第一の積分で $\sin t$ を微分し，$\e^t$ を積分すると，
\[
\begin{aligned}
 I(x)
 &=[\e^t\sin t]_0^x-\int_0^x\e^t\cos t\dd t\\
 &=\e^x\sin x-J(x).
\end{aligned}
\]
続いて $J(x)$ でも $\e^t$ を積分する側に選ぶ。$\cos t$ の微分は $-\sin t$ だから，
\[
\begin{aligned}
 J(x)
 &=[\e^t\cos t]_0^x-\int_0^x\e^t(-\sin t)\dd t\\
 &=\e^x\cos x-1+I(x).
\end{aligned}
\]
ここで $I(x)$ が戻ってきた。第二式を第一式へ代入すると
\[
\begin{aligned}
 I(x)
 &=\e^x\sin x-\{\e^x\cos x-1+I(x)\},\\
 2I(x)&=\e^x(\sin x-\cos x)+1.
\end{aligned}
\]
よって $I(x)=\frac12\e^x(\sin x-\cos x)+\frac12$ である。最後の $\frac12$ も任意定数へまとめれば，
\[
 \boxed{\int\e^x\sin x\dd x
 =\frac12\e^x(\sin x-\cos x)+C}
\]
を得る。微分による確認は
\[
 \frac12\e^x(\sin x-\cos x)
 +\frac12\e^x(\cos x+\sin x)
 =\e^x\sin x
\]
である。

ここで役立ったのは，「同じ積分が戻る」という事実だけではない。右辺へ $-I(x)$ として戻ったので，移項すると $2I(x)$ が得られたことが重要である。もし同じ式を単に逆向きにたどり，$I(x)=I(x)$ に戻っただけなら，新しい情報はない。循環した後には，何倍の $I(x)$ が残ったのかを確かめる。

不定積分のまま同じ計算を進めることもできる。その場合は，途中の等式が任意定数の違いを許していることを忘れず，最後に一つの $C$ にまとめる。定積分として $I,J$ を定めたのは，等式の両辺を同じ関数として固定し，移項の意味をはっきりさせるためである。

\section{置換で形を変えてから部分積分する}\label{節:置換と部分積分}

積分法は，問題ごとに一つだけ選ぶものではない。置換によって合成関数をほどくと，今度は二つの関数の積が現れ，部分積分が使えることがある。前章でも，根号を変数に換えた後で割り算を使った。同じように，途中の式に合わせて次の道具を選べばよい。

\begin{example}{対数を角に持つ三角関数}
$x>0$ において，$\int\sin(\log x)\dd x$ を求めよ。
\end{example}

角の $\log x$ を一つの変数 $t$ にしたい。しかし，もとの積分には $\frac1x$ がないので，$\sin t\dd t$ と置き換えることはできない。$t=\log x$ を $x=\e^t$ と解き直すと，
\[
 \dd x=\e^t\dd t
\]
である。この微分の係数も含めて換えると，
\[
 \int\sin(\log x)\dd x
 =\int\e^t\sin t\dd t
\]
となる。これは直前の例題で求めた形だから，
\[
\begin{aligned}
 \int\sin(\log x)\dd x
 &=\frac12\e^t(\sin t-\cos t)+C\\
 &=\frac x2\{\sin(\log x)-\cos(\log x)\}+C.
\end{aligned}
\]
$\log x$ は $x>0$ で狭義単調増加し，逆向きの対応 $x=\e^t$ も一つに決まっている。置換後に $\e^t$ が現れたのは，新しい問題を余計に作ったためではなく，$\dd x$ を正しく書き換えたためである。

検算では，外にある $x$ と，内側の $\log x$ の両方を微分する。
\[
\begin{aligned}
 &\frac{\dd}{\dd x}
 \left[\frac x2\{\sin(\log x)-\cos(\log x)\}\right]\\
 &\qquad=\frac12\{\sin(\log x)-\cos(\log x)\}
   +\frac x2\cdot\frac1x
     \{\cos(\log x)+\sin(\log x)\}\\
 &\qquad=\sin(\log x).
\end{aligned}
\]
$t=\log x$ と置くだけでは終わらず，置換後の積を部分積分で求め，最後に元の変数へ戻した。この順序が，それぞれの手法の役割を分けている。

\section{根号を含む積分を同じ積分へ戻す}\label{節:平方根と部分積分}

$\sqrt{1+x^2}$ も，そのままでは基本公式の形ではない。前章の三角置換から $x=\tan\theta$ を考えることはできるが，ここでは部分積分を使う別の道を調べる。手がかりは，根号を微分すると $\frac{x}{\sqrt{1+x^2}}$ になり，さらに $x$ を掛ければ分子に $x^2$ ができることである。この $x^2$ を $(1+x^2)-1$ と分けると，もとの根号が再び現れる。

\begin{example}{根号とその逆数を関係付ける}
$\int\sqrt{1+x^2}\dd x$ を求めよ。
\end{example}

まず $u=\sqrt{1+x^2}$，$v'=1$ と選ぶ。$u'=\frac{x}{\sqrt{1+x^2}}$，$v=x$ だから，部分積分から
\[
 \int\sqrt{1+x^2}\dd x
 =x\sqrt{1+x^2}-\int\frac{x^2}{\sqrt{1+x^2}}\dd x
\]
を得る。右辺の積分を見ただけでは簡単になったか分かりにくいが，分子を分けると
\[
 \frac{x^2}{\sqrt{1+x^2}}
 =\frac{(1+x^2)-1}{\sqrt{1+x^2}}
 =\sqrt{1+x^2}-\frac1{\sqrt{1+x^2}}
\]
となる。したがって，もとの積分を移項すれば，任意定数を除いて
\[
 2\int\sqrt{1+x^2}\dd x
 =x\sqrt{1+x^2}+\int\frac1{\sqrt{1+x^2}}\dd x
\]
という関係が得られる。未知の積分がもう一度現れたが，その係数をまとめることができた。残る課題は，根号の逆数の積分である。

ここで $r(x)=\sqrt{1+x^2}$ と略記して，$x+r(x)$ を微分してみる。
\[
 \{x+r(x)\}'=1+\frac{x}{r(x)}
 =\frac{x+r(x)}{r(x)}.
\]
この形には，微分する前の $x+r(x)$ が分子に現れている。そこで対数の微分を使えば，その因子が約分されて
\[
 \frac{\dd}{\dd x}\log\{x+r(x)\}
 =\frac{1}{x+r(x)}\cdot\frac{x+r(x)}{r(x)}
 =\frac1{r(x)}
\]
となる。しかも $r(x)=\sqrt{1+x^2}>|x|$ なので，$x+r(x)>0$ はすべての実数 $x$ で成り立つ。対数の中身の符号も確定した。よって
\[
 \int\frac1{\sqrt{1+x^2}}\dd x
 =\log\left(x+\sqrt{1+x^2}\right)+C.
\]
これを先ほどの関係へ戻すと，
\[
 \boxed{
 \int\sqrt{1+x^2}\dd x
 =\frac12\left\{x\sqrt{1+x^2}
       +\log\left(x+\sqrt{1+x^2}\right)\right\}+C
 }
\]
を得る。

$x+r(x)$ は，突然覚えるべき式として出てきたのではない。$r'(x)=\frac{x}{r(x)}$ と $x'=1=\frac{r(x)}{r(x)}$ を足すと，分子がもとの和になる。そのため，対数の微分に現れる「もとの式で割る」という操作が，求めたい $\frac1{r(x)}$ を作るのである。

最後に答え全体を微分すると，
\[
\begin{aligned}
 &\frac12\left\{\sqrt{1+x^2}
       +\frac{x^2}{\sqrt{1+x^2}}
       +\frac1{\sqrt{1+x^2}}\right\}\\
 &\qquad=\frac12\{\sqrt{1+x^2}+\sqrt{1+x^2}\}
 =\sqrt{1+x^2}
\end{aligned}
\]
となる。この確認により，途中で不定積分を移項した計算も，最終的な原始関数として正しいことが分かる。後で双曲線関数を学ぶと，同じ根号を別の置換で扱える。ここで得た式は，その別解と比較する基準になる。

\section{定積分では端点の項を先に見る}\label{節:定積分の部分積分}

不定積分を求めてから上下端を代入してもよいが，最初から定積分の公式を使えば，端点で消える項がすぐに分かる。公式
\[
 \int_a^b uv'\dd x=[uv]_a^b-\int_a^b u'v\dd x
\]
では，残る積分だけでなく，端点の $u(a),u(b),v(a),v(b)$ も計算の見通しに関わる。積のどちらかが $0$ ならば，その端点の寄与はなくなる。

\begin{example}{端点で消える積を確かめる}
$\int_0^{\frac\pi2}x\sin x\dd x$ を求めよ。
\end{example}

$u=x$，$v'=\sin x$ と選ぶと，$u'=1$，$v=-\cos x$ なので，
\[
\begin{aligned}
 \int_0^{\frac\pi2}x\sin x\dd x
 &=[-x\cos x]_0^{\frac\pi2}
   +\int_0^{\frac\pi2}\cos x\dd x\\
 &=0+[\sin x]_0^{\frac\pi2}=1.
\end{aligned}
\]
上端では $\cos\frac\pi2=0$，下端では $x=0$ のため，端点の項はともに $0$ になった。二つの端点で消える理由は異なるので，一方だけを見て項全体を消しているわけではない。

また，積分される関数はこの区間で非負だから，求めた値が正であることは自然である。仮に $-1$ になったら，$\sin x$ の原始関数の負号か，部分積分の前の負号を確かめるべきだと分かる。

\begin{example}{部分積分の後に割り算を使う}
$\int_0^1x\log(1+x)\dd x$ を求めよ。
\end{example}

$\log(1+x)$ を微分すると $\frac1{1+x}$ になる。$x$ の原始関数は $\frac{x^2}{2}$ だから，残る積分は有理関数になる。そこで
\[
 u=\log(1+x),\qquad v'=x
\]
と選び，
\[
 \int_0^1x\log(1+x)\dd x
 =\left[\frac{x^2}{2}\log(1+x)\right]_0^1
  -\frac12\int_0^1\frac{x^2}{1+x}\dd x
\]
を得る。第一項は $\frac12\log2$ である。第二項では，分子の次数が分母より高いので，\chapterref{章:基本積分}で学んだ割り算を使う。
\[
 x^2=(1+x)(x-1)+1,
 \qquad
 \frac{x^2}{1+x}=x-1+\frac1{1+x}.
\]
したがって
\[
\begin{aligned}
 \int_0^1\frac{x^2}{1+x}\dd x
 &=\left[\frac{x^2}{2}-x+\log(1+x)\right]_0^1\\
 &=-\frac12+\log2.
\end{aligned}
\]
これを戻すと
\[
 \int_0^1x\log(1+x)\dd x
 =\frac12\log2-\frac12\left(-\frac12+\log2\right)
 =\frac14.
\]
部分積分の直後には対数が端点に残り，別の積分にも対数が現れたが，最後に差を取ると打ち消し合った。途中の式の見た目だけで，答えにも対数が必ず残ると決めることはできない。

ここまでで，式の整理，置換積分，部分積分という三つの道具を学んだ。整式の割り算や三角恒等式で和へ直せるなら，まずその変形を試せる。合成関数と内側の微分がそろっていれば置換を考え，微分すると簡単になる因子と，積分できる因子の積なら部分積分を考える。ある操作で何が簡単になり，何が残るかを具体的に見ることが，選択の根拠になる。

ただし，定積分の値だけが欲しいときは，原始関数を最後まで求める必要がない場合もある。区間の左右を入れ換えたときに式がどう変わるか，正負の部分がどう対応するかを使えば，値が決まることがある。次の\chapterref{章:定積分の対称性}では，被積分関数と積分区間を一緒に見て，その対応を計算へ生かす。

\section*{章末問題}
\begin{enumerate}[label=\textbf{\thechapter.\arabic*},leftmargin=*]
 \item\label{問:部分一次指数} $\int x\e^{3x}\dd x$ を求めよ。微分する因子と積分する因子を示し，答えを微分して確かめよ。
 \item\label{問:部分整式対数} $x>0$ において，$\int x^2\log x\dd x$ を求めよ。$x^2$ と $\log x$ のどちらを微分する側に選ぶか，その理由も述べよ。
 \item\label{問:部分反復} $\int x^3\sin x\dd x$ を求めよ。部分積分を繰り返すたびに，整式の次数がどう変わるかを確かめよ。
 \item\label{問:部分循環} 部分積分を二回使い，$\int\e^{2x}\cos x\dd x$ を求めよ。もとの積分が戻った後の等式も書け。
 \item\label{問:部分対数余弦} $x>0$ において，$\int\cos(\log x)\dd x$ を求めよ。置換の際には，$\dd x$ が何に変わるかも書け。
 \item\label{問:部分置換定積分} $\int_0^1x\log(1+x^2)\dd x$ を求めよ。置換の後に，どの積分を部分積分で求めればよいかを示せ。
 \item\label{問:部分根号定積分} 部分積分を使い，$\int_0^1\sqrt{4+x^2}\dd x$ を求めよ。必要になる根号の逆数の原始関数は，微分によって確かめて用いよ。
 \item\label{問:部分情報だけの計算} $f$ は $[0,1]$ で連続微分可能とし，$f(0)=1$，$f(1)=2$，$\int_0^1f(x)\dd x=\frac32$ を満たす。$f$ の具体的な式を求めずに，$\int_0^1xf'(x)\dd x$ と $\int_0^1(1-x)f'(x)\dd x$ を求めよ。
\end{enumerate}

\smallskip\noindent 解答は\pageref{解答:第6章}ページ。
